\hypertarget{appendix-a-tex-internals-for-programmers}{%
\chapter{TeX Internals for Programmers}\label{appendix-a-tex-internals-for-programmers}}

Chapter 1 introduced boxes, glue, penalties, tokens, and category codes at the level needed to start writing and debugging LaTeX documents productively. This appendix goes further, into the actual mechanics underneath that introduction, for the reader who wants to understand not just that a particular behavior happens but exactly why, down to the primitive level. It is written for a reader comfortable with programming language internals generally, tokenizers, parsers, evaluators, and treats TeX as what it is: a small, unusually literal, and unusually stable programming language whose only unusual feature is that its typical output is a typeset page rather than a return value.

Nothing in the main chapters requires this appendix. It exists for Chapters 1, 6, 8, and 19 specifically, for a reader who wants the full mechanism behind the model those chapters use.

\hypertarget{a.1-the-category-code-table-in-full}{%
\section{The Category Code Table in Full}\label{a.1-the-category-code-table-in-full}}

Every character TeX reads is assigned one of sixteen category codes at the moment it is read, determined by a table that can be modified at any point during processing. The categories are: 0, escape character, marking the start of a control sequence; 1, begin group, ordinarily \texttt{\{}; 2, end group, ordinarily \texttt{\}}; 3, math shift, ordinarily \texttt{\$}; 4, alignment tab, ordinarily \texttt{\&}; 5, end of line; 6, parameter, ordinarily \texttt{\#}; 7, superscript, ordinarily \texttt{\^{}}; 8, subscript, ordinarily \texttt{\_}; 9, ignored; 10, space; 11, letter; 12, other; 13, active, a character that behaves as a control sequence in its own right; 14, comment character, ordinarily \texttt{\%}; 15, invalid character; and a small number of characters TeX treats specially at the input-line level before category codes are even consulted.

Two consequences follow directly from this table's mutability. First, \texttt{\textbackslash{}catcode} assignments are what packages like \texttt{verbatim} actually rely on: temporarily reassigning most characters to category 12 (other) suspends their special meaning entirely, which is why text inside a \texttt{verbatim} environment can contain unescaped \texttt{\textbackslash{}}, \texttt{\{}, \texttt{\}}, and \texttt{\%} without any of them being interpreted. Second, active characters, category 13, are the mechanism behind packages that make a character like \texttt{\textasciitilde{}} or \texttt{"} behave as a macro invocation, \texttt{babel}'s language-specific shortcuts and \texttt{csquotes}'s smart quote handling both work by making an otherwise ordinary character active and defining what it expands to.

\hypertarget{a.2-tokenization}{%
\section{Tokenization}\label{a.2-tokenization}}

TeX's input processor reads one line at a time, applies the category code table current at that moment to convert the line's characters into tokens, and hands the resulting token list to the expansion and execution machinery described in the following sections. A token is either a control sequence, formed by an escape character followed by a run of letter-category characters (or a single non-letter character), or a character token, a pairing of a character and its category code.

This is worth stating precisely because it clarifies a common point of confusion: \texttt{\textbackslash{}ab} and \texttt{\textbackslash{}a} followed by \texttt{b} are not automatically the same two-token or one-token sequence; whether \texttt{\textbackslash{}ab} tokenizes as the single control sequence \texttt{\textbackslash{}ab} or as \texttt{\textbackslash{}a} followed by a literal \texttt{b} depends entirely on whether \texttt{b} has category 11 (letter) at the moment that line is read, since control sequence names absorb a run of consecutive letter-category characters. This is also why an explicit space or a non-letter character is needed after certain control sequences to prevent TeX from trying to absorb further letters into the control sequence name that was not intended, and why \texttt{\textbackslash{}} (control sequence followed by an explicit space) is TeX's idiom for ``a control sequence name ends here, and additionally insert a space,'' distinct from simply writing the next character with nothing between.

\hypertarget{a.3-expansion-in-detail}{%
\section{Expansion in Detail}\label{a.3-expansion-in-detail}}

Chapter 6 introduced expansion as TeX's core mechanism: an expandable token is replaced by what it expands to, repeatedly, until nothing further expandable remains at the position currently under consideration. The detail worth adding here is exactly which tokens count as expandable, since this determines what \texttt{\textbackslash{}edef}, \texttt{\textbackslash{}expandafter}, and the various \texttt{\textbackslash{}if} conditionals can and cannot reach.

Expandable tokens include macros (anything defined with \texttt{\textbackslash{}def} or built on top of it), the conditional primitives (\texttt{\textbackslash{}ifnum}, \texttt{\textbackslash{}ifdim}, \texttt{\textbackslash{}ifx}, and the rest of the \texttt{\textbackslash{}if} family), \texttt{\textbackslash{}csname...\textbackslash{}endcsname} for constructing a control sequence name from its component characters, \texttt{\textbackslash{}expandafter} itself, and a handful of others such as \texttt{\textbackslash{}input} (which expands to the contents of the named file) and \texttt{\textbackslash{}the} (which expands to the printable representation of a register's current value). Non-expandable tokens include ordinary character tokens once they have reached this stage, box-construction and assignment primitives, and anything that has a side effect on TeX's internal state rather than producing further tokens.

This distinction is what makes \texttt{\textbackslash{}expandafter} meaningful rather than merely a curiosity: \texttt{\textbackslash{}expandafter\textbackslash{}A\textbackslash{}B} expands \texttt{\textbackslash{}B} once (if it is expandable) before \texttt{\textbackslash{}A} is examined at all, letting a macro writer reach past one token to affect what comes immediately after it, which is the standard technique for constructing a control sequence name out of pieces via \texttt{\textbackslash{}csname} or for testing the category or value of whatever comes next in the input before committing to how the current macro should proceed. Chained \texttt{\textbackslash{}expandafter} calls, \texttt{\textbackslash{}expandafter\textbackslash{}expandafter\textbackslash{}expandafter\textbackslash{}A\textbackslash{}expandafter\textbackslash{}B\textbackslash{}C}, extend the same idea to reach further into an input stream than a single application would allow, at the cost of a syntax that is notoriously difficult to read without deliberate practice; the \texttt{expl3} function-based interface exists largely to make constructions like this unnecessary in ordinary macro-writing.

\hypertarget{a.4-grouping-and-scope}{%
\section{Grouping and Scope}\label{a.4-grouping-and-scope}}

\texttt{\{} and \texttt{\}}, category codes 1 and 2, delimit a group, and a group is TeX's only mechanism for scoping an assignment. Any assignment made inside a group, a font change, a counter change, a redefinition of an existing macro via local \texttt{\textbackslash{}def}, is automatically undone at the group's closing brace, restored to whatever value it held before the group opened. This is why \texttt{\{\textbackslash{}bfseries\ bold\ text\}\ normal\ text} correctly returns to the surrounding font without any explicit ``turn off bold'' command: the closing brace itself performs the restoration, because the font change was a local assignment scoped to that group.

\texttt{\textbackslash{}global} overrides this behavior for a specific assignment, making it persist past the end of the current group rather than being undone; \texttt{\textbackslash{}def} versus \texttt{\textbackslash{}gdef}, and the various registers' local versus global assignment forms, follow the same pattern throughout TeX. Environments (\texttt{\textbackslash{}begin\{...\}...\textbackslash{}end\{...\}}) are, at the primitive level, built from exactly this grouping mechanism, \texttt{\textbackslash{}begin\{env\}} opens a group and \texttt{\textbackslash{}end\{env\}} closes it, which is why an assignment made inside an environment is automatically local to it without any special environment-specific scoping machinery being required.

\hypertarget{a.5-registers-texs-only-variables}{%
\section{Registers: TeX's Only Variables}\label{a.5-registers-texs-only-variables}}

TeX has no general-purpose variables in the sense a conventional programming language would; it has a fixed set of typed register classes, each with a fixed number of numbered slots, and LaTeX's counters, lengths, and named boxes are all built on top of these registers rather than being a separate mechanism. Count registers hold integers and back LaTeX's \texttt{\textbackslash{}newcounter} and every \texttt{\textbackslash{}value}-accessible counter. Dimen registers hold a length with an attached unit and back \texttt{\textbackslash{}newlength}. Skip registers hold glue, a dimension plus stretch and shrink components, and back the vertical and horizontal spacing lengths a document defines. Box registers hold a constructed box, and are what \texttt{\textbackslash{}newsavebox} and \texttt{\textbackslash{}sbox}/\texttt{\textbackslash{}savebox} allocate and populate. Toks registers hold a token list, unexpanded, and are used less often directly by document authors but underlie some of the more advanced macro-programming techniques in Part III.

\texttt{\textbackslash{}newcount}, \texttt{\textbackslash{}newdimen}, and their relatives, which LaTeX's \texttt{\textbackslash{}newcounter} and \texttt{\textbackslash{}newlength} are thin wrappers around, allocate a previously unused numbered register and bind a control sequence name to it, which is why TeX historically imposed a hard limit on the total number of registers of each type available at once (256 in classic implementations, considerably larger in modern engines including LuaTeX), and why careless repeated allocation of registers inside a macro that gets called many times, rather than allocating once at load time, was historically a real resource concern, though less pressing in current engines than it was in the era the limit was originally set.

\hypertarget{a.6-dimensions-and-glue}{%
\section{Dimensions and Glue}\label{a.6-dimensions-and-glue}}

A TeX dimension is a number followed by a unit: \texttt{pt} (point), \texttt{in}, \texttt{cm}, \texttt{mm}, \texttt{em} and \texttt{ex} (relative to the current font's size), among others, all ultimately converted internally to a common fixed-point representation in units of scaled points. Arithmetic on dimensions, adding two lengths, multiplying a length by an integer, is supported directly by TeX's register and assignment syntax, though it is limited compared to a general-purpose numeric type; \texttt{calc} and, in the LaTeX3 kernel, the expandable arithmetic in \texttt{l3fp}, extend this considerably for macro-writing purposes.

Glue, introduced conceptually in Chapter 1, is formally a dimension with two additional components: a stretch amount and a shrink amount, each itself a dimension, optionally given in ``fil'' units (\texttt{fil}, \texttt{fill}, \texttt{filll}) representing infinite stretch or shrink of increasing order, used for constructions like \texttt{\textbackslash{}hfill} that should absorb all remaining space on a line regardless of how much space that turns out to be. When TeX justifies a line or a page, it computes the total natural width of the material on that line or page, compares it to the target width, and distributes the required stretch or shrink proportionally across all the glue present, weighted by each glue item's stretch or shrink component, which is the actual mechanism behind the paragraph justification described at a higher level in Chapter 1.

\hypertarget{a.7-penalties-and-the-line-breaking-algorithm}{%
\section{Penalties and the Line-Breaking Algorithm}\label{a.7-penalties-and-the-line-breaking-algorithm}}

TeX's paragraph-breaking algorithm, generally credited to Knuth and Plass~\cite{knuthplass1981}, does not choose line breaks greedily, filling each line as full as possible and moving to the next; it considers the entire paragraph at once, evaluating the total ``badness'' of every combination of break points and choosing the set of breaks that minimizes an accumulated cost function across the whole paragraph, subject to the penalties attached to each candidate break point. A break with an attached penalty near the forcing threshold (\(-10000\) or beyond) is effectively mandatory; a large positive penalty makes a break progressively less likely to be chosen even where the surrounding glue could accommodate it; penalties near zero leave the decision almost entirely to the badness calculation.

This global, whole-paragraph optimization is why TeX's line breaking is generally regarded as superior to a simpler greedy algorithm: a greedy algorithm can leave one line very loose and the next disproportionately tight without any way to trade slack between them, while Knuth-Plass considers the paragraph as a whole and can accept slightly more stretch on one line in exchange for a much better break somewhere else in the same paragraph. The page-breaking algorithm applies a related but distinct badness-and-penalty evaluation to the vertical dimension, considering where to break across pages given the accumulated material and any explicit penalties (widow and orphan control, \texttt{\textbackslash{}nopagebreak}, and so on) that discourage or force a page break at a specific point.

\hypertarget{a.8-boxes-and-box-construction}{%
\section{Boxes and Box Construction}\label{a.8-boxes-and-box-construction}}

A box, as introduced in Chapter 1, is either an \texttt{\textbackslash{}hbox} (horizontal box, built by placing its contents side by side) or a \texttt{\textbackslash{}vbox} (vertical box, built by stacking its contents), each with a width, height, and depth (the depth being the extent below the baseline, relevant for characters like \texttt{g} or \texttt{y} that descend below it). \texttt{\textbackslash{}vtop} is a variant of vertical box construction whose reference point is its first line rather than its last, useful when aligning the tops rather than the bottoms of adjacent material.

Every piece of typeset material is ultimately assembled into boxes recursively: a character becomes a small box supplied by the current font, a word is an \texttt{\textbackslash{}hbox} of character boxes with interword glue, a line is an \texttt{\textbackslash{}hbox} of word boxes, a paragraph is a \texttt{\textbackslash{}vbox} of line boxes, and a page is, at the top level, a \texttt{\textbackslash{}vbox} containing everything on it. \texttt{\textbackslash{}raise}, \texttt{\textbackslash{}lower}, \texttt{\textbackslash{}moveleft}, and \texttt{\textbackslash{}moveright} shift an already-constructed box relative to its surrounding context without changing its own internal structure, and are the primitives underlying LaTeX-level positioning commands used for fine typographic adjustment.

\hypertarget{a.9-the-page-builder-and-output-routine}{%
\section{The Page Builder and Output Routine}\label{a.9-the-page-builder-and-output-routine}}

Material accumulates in what TeX calls the ``main vertical list'' as a document is processed, and the page builder periodically checks whether enough material has accumulated to fill a page, using the same badness-and-penalty evaluation described in A.7 applied to page-length material rather than paragraph-length material. When it decides a page break should occur, it invokes the output routine, a token list, customizable via \texttt{\textbackslash{}output}, responsible for taking the box representing everything accumulated for that page (available inside the output routine as \texttt{\textbackslash{}box255}) and actually placing it, adding headers, footers, and page numbers, and sending the result to the DVI or PDF driver.

This is the mechanism packages that modify page layout, custom headers and footers, multi-column output, watermarking, actually hook into: \texttt{fancyhdr} and similar packages redefine \texttt{\textbackslash{}output} or the hooks it calls, rather than modifying anything about how individual pages are composed at a lower level. Understanding that headers and footers are inserted by the output routine, at the point a page is finalized, rather than being part of each page's content from the start, explains why certain header/footer customizations interact awkwardly with certain float placements or multi-column layouts: both are competing for control over material that has already been provisionally assigned to a specific page by the page builder, and the output routine is the single point where that assignment can still be revised.

\hypertarget{a.10-a-first-script-in-luatex}{%
\section{A First Script in LuaTeX}\label{a.10-a-first-script-in-luatex}}

LuaTeX, the engine underlying LuaLaTeX, embeds a full Lua interpreter and exposes hooks into nearly every stage of TeX's processing described in this appendix: the input processor, the tokenizer, the paragraph and page builders, and the output routine can all be observed or modified from Lua code, in addition to Lua being usable simply as a general-purpose scripting language callable from within a document via \texttt{\textbackslash{}directlua}.

A minimal first use of this capability, connecting directly to the build-time computation described in Chapters 14 and 15, is reading an external value into the document from Lua rather than through LaTeX's own, more limited macro-programming facilities: \texttt{\textbackslash{}directlua\{tex.sprint(2\^{}10)\}} evaluates a Lua expression and inserts its printed result directly into the document, which for anything beyond simple arithmetic, reading a JSON or CSV file's contents directly rather than through a dedicated LaTeX package, calling out to an external library available to Lua, becomes a genuinely useful capability rather than a curiosity. This is, deliberately, as far as this appendix goes into LuaTeX scripting; a reader who wants to build substantially on this capability is better served by the \texttt{luatex-plain} and LuaTeX-specific documentation than by an appendix whose purpose is to explain the primitives underneath ordinary LaTeX use, not to teach Lua programming.

With this appendix's model of tokens, expansion, grouping, registers, boxes, and the page builder in hand, the behavior described throughout Chapters 1 through 19, why a macro expands the way it does, why a fragile command breaks under reprocessing, why a debugging session traces expansion rather than execution, should read as the direct, predictable consequence of a small number of rules applied consistently, rather than as a collection of separately memorized special cases.
